% !TEX program = pdflatex
% !TEX encoding = UTF-8 Unicode

% Plantilla de la clase `scrbook` del paquete KOMA-script para la
% elaboración de un TFG siguiendo las directrices del la comisión del
% Grado en Matemáticas de la Universidad de Granada.

% Francisco Torralbo Torralbo
% miércoles, 29 de abril de 2020

\documentclass[twoside=semi]{scrbook}

\KOMAoptions{%
  fontsize=10pt,        % Tamaño de fuente
  paper=a4,             % Tamaño del papel
  headings=normal,      % Tamaño de letra para los títulos: small, normal, big
  parskip=half,         % Espacio entre párrafos: full (una línea) o half (media línea)
  headsepline=false,    % Una linea separa la cabecera del texto
  cleardoublepage=empty,% No imprime cabecera ni pie en páginas en blanco 
  chapterprefix=false,  % No antepone el texto "capítulo" antes del número
  appendixprefix=false,	% No antepone el texto "Apéndice" antes de la letra
  listof=totoc,		    	% Añade a la tabla de contenidos la lista de tablas y figuras
  index=totoc,			    % Añade a la tabla de contenidos una entrada para el índice
  bibliography=totoc,	  % Añade a la tabla de contenidos una entrada para bibliografía
  BCOR=0mm,           % Reserva margen interior para la encuadernación. 
                        % El valor dependerá el tipo de encuadernado y del grosor del libro.
  DIV=10,             % Cálcula el diseño de página según ciertos 
                        % parámetros. Al aumentar el número aumentamos el ancho de texto y disminuimos el ancho del margen. Una opción de 14 producirá márgenes estrechos y texto ancho.
}

% INFORMACIÓN PARA LA VERSIÓN IMPRESA
% Si el documento ha de ser impreso en papel de tamaño a4 pero el tamaño del documento (elegido en \KOMAoptions con la ocpión paper) no es a4 descomentar la línea que carga el paquete `crop` más abajo. El paquete crop se encargará de centrar el documento en un a4 e imprimir unas guías de corte. El procedimiento completo para imprenta sería el siguiente:
% 0. Determinar, según el tipo de encuadernación del documento, el ancho reservado para el proceso de encuadernación (preguntar en la imprenta), es decir, la anchura del área del papel que se pierde durante el proceso de encuadernación. Fijar la varibale BCOR de \KOMAoptions a dicho valor.
% 1. Descomentar la siguiente línea e imprimir una única página con las guías de corte
% 2. Cambiar la opción `cross` por `cam` (o `off`) en el paquete crop y volver a compilar. Imprimir el documento (las guías de corte impresas no inferfieren con el texto).
% 3. Usar la página con las guías impresas en el punto 1 para cortar todas las páginas.

%\usepackage[a4, odd, center, pdflatex, cross]{crop} % Permite imprimir el documento en un a4 (si el tamaño es más pequeño) mostrando unas guías de corte. Útil para imprenta.

% VERSIÓN ELECTRÓNICA PARA TABLETA
% Las opciones siguientes seleccionan un tamaño de impresión similar a una tableta de 9 pulgadas con márgenes estrechos. Útil para producir una versión en pdf para ser leída en una tableta en lugar de impresa.
% Para que la portada quede centrada correctamente hay que editar el
% archivo `portada.tex` y eliminar el entorno `addmargin`

% \KOMAoptions{fontsize=10pt, paper=19.7104cm:14.7828cm, twoside=false, BCOR=0cm, DIV=14}

% ---------------------------------------------------------------------
%	PAQUETES 
% ---------------------------------------------------------------------

% CODIFICACIÓN E IDIOMA
% ---------------------------------------------------------------------
\usepackage[utf8]{inputenc} 			    % Codificación de caracteres

% Selección del idioma: cargamos por defecto inglés y español (aunque este último es el idioma por defecto para el documento). Cuando queramos cambiar de idioma escribiremos:
% \selectlanguage{english} o \selectlanguage{spanish}

\usepackage[english, spanish, es-nodecimaldot, es-noindentfirst, es-tabla]{babel}

% Opciones cargadas para el paquete babel:
  % es-nodecimaldot: No cambia el punto decimal por una coma en modo matemático.
  % es-noindentfirst: No sangra los párrafos tras los títulos.
  % es-tabla: cambia el título del entorno `table` de "Cuadro" a "Tabla"

% Otras opciones del paquete spanish-babel:
  \unaccentedoperators % Desactiva los acentos en los operadores matemáticso (p.e. \lim, \max, ...). Eliminar esta opción si queremos que vayan acentuados

% MATEMÁTICAS
% ---------------------------------------------------------------------
\usepackage{amsmath, amsthm, amssymb} % Paquetes matemáticas
\usepackage{mathtools}                % Añade mejoras a amsmath
\mathtoolsset{showonlyrefs=true}      % sólo se numeran las ecuaciones que se usan
\usepackage[mathscr]{eucal} 					% Proporciona el comando \mathscr para
                                      % fuentes de tipo manuscrito en modo matemático sin sobreescribir el comando \mathcal

% TIPOGRAFÍA 
% ---------------------------------------------------------------------
% El paquete microtype mejora la tipografía del documento.
\usepackage[activate={true,nocompatibility},final,tracking=true,kerning=true,spacing=true,factor=1100,stretch=10,shrink=10]{microtype}

% Las tipografías elegidas para el documento, siguiendo la guía de estilo de la UGR,
% son las siguientes
% Normal font: 			URW Palladio typeface. 
% Sans-serif font: 	Gill Sans
% Monospace font: 	Inconsolata
\usepackage[T1]{fontenc}
\usepackage[sc, osf]{mathpazo} \linespread{1.05}         
\usepackage[scaled=.95,type1]{cabin} % sans serif in style of Gill Sans
% Si el paquete cabin da error usar el siguiente comando en su lugar
% \renewcommand{\sfdefault}{iwona} 
\usepackage{inconsolata}


% Selecciona el tipo de fuente para los títulos (capítulo, sección, subsección) del documento.
\setkomafont{disposition}{\sffamily\bfseries}

% Cambia el ancho de la cita. Al inicio de un capítulo podemos usar el comando \dictum[autor]{cita} para añadir una cita famosa de un autor.
\renewcommand{\dictumwidth}{0.45\textwidth} 

\recalctypearea % Necesario tras definir la tipografía a usar.

\usepackage{setspace}

% OPCIONES AVANZADAS DE ENUMERATE
\usepackage{enumerate}

% TABLAS, GRÁFICOS Y LISTADOS DE CÓDIGO
% ---------------------------------------------------------------------
\usepackage{booktabs}
% \renewcommand{\arraystretch}{1.5} % Aumenta el espacio vertical entre las filas de un entorno tabular

\usepackage{xcolor, graphicx}
% Carpeta donde buscar los archivos de imagen por defecto
\graphicspath{{img/}}

% Para poder posiciones imagenes y tablas con H (HERE)
\usepackage{float}

% Para tablas con \hline
\usepackage{hhline}

% Para utilizar el entorno comment
\usepackage{verbatim}

% Para insertar un pdf dentro del documento
\usepackage{pdfpages}

% IMAGEN DE LA PORTADA
% Existen varias opciones para la imagen de fondo de la portada del TFG. Todas ellas tienen en logotipo de la universidad de Granada en la cabecera. Las opciones son las siguientes:
% 1. portada-ugr y portada-ugr-color: diseño con marca de agua basada en el logo de la UGR (en escala de grises y color).
% 2. portada-ugr-sencilla y portada-ugr-sencilla-color: portada únicamente con el logotipo de la UGR en la cabecera.
\usepackage{eso-pic}
\newcommand\BackgroundPic{%
	\put(0,0){%
		\parbox[b][\paperheight]{\paperwidth}{%
			\vfill
			\centering
      % Indicar la imagen de fondo en el siguiente comando
			\includegraphics[width=\paperwidth,height=\paperheight,%
			keepaspectratio]{portada-ugr-color}%
			\vfill
}}}

\usepackage{listings} % Para la inclusión de trozos de código

% CABECERAS
% ---------------------------------------------------------------------
% Si queremos modificar las cabeceras del documento podemos usar el paquete
% `scrlayer-scrpage` de KOMA-Script. Consultar la documentación al respecto.
% \usepackage[automark]{scrlayer-scrpage}

% VARIOS
% ---------------------------------------------------------------------

%\usepackage{showkeys}	% Muestra las etiquetas del documento. Útil para revisar las referencias cruzadas.

% ÍNDICE 
% Para generar el índice hay que compilar el documento con MakeIndex. Generalmente los editores se encargan de ello automáticamente.
% ----------------------------------------------------------------------
% \index{} para añadir un elemento
% \index{main!sub} para añadir un elementos "sub" bajo la categoría "main".
% \index{termino|textbf} para dar formato al número de página (negrita).
% \index{termino|see{termino relacionado}} para crear una referencia cruzada

% Ejemplo: \index{espacio homogéneo}, \index{superficie!mínima}, \index{esfera|see{espacio homogéneo}}

% Activar los siguientes comandos para generar el índice terminológico. Ver también comandos al final de este documento para incluir dicho índice en el pdf final.
\usepackage{makeidx}
\makeindex

% Para revisar las entradas al índice conforme las incluimos en el documento es útil el siguiente paquete. Conviene observar que mientras esté cargado no se generará el índice.
%\usepackage{showidx} % Muestra en el margen del documento las entradas añadidas al índice. Útil para revisar el documento. Si está activo el índice no se genera

% ---------------------------------------------------------------------
% COMANDOS Y ENTORNOS
% ---------------------------------------------------------------------
% Cargamos un archivo externo donde hemos incluido todos los comandos
% propios que vamos a usar en el documento.
% DEFINICIÓN DE COMANDOS Y ENTORNOS

% CONJUNTOS DE NÚMEROS
  \newcommand{\N}{\mathbb{N}}     % Naturales
  \newcommand{\R}{\mathbb{R}}     % Reales
  \newcommand{\Z}{\mathbb{Z}}     % Enteros
  \newcommand{\Q}{\mathbb{Q}}     % Racionales
  \newcommand{\C}{\mathbb{C}}     % Complejos

% Otros comandos propios
  \newcommand{\mun}{\{\mu_n\}_{n \in \N_0}}                  				        % Sucesión de momentos
  \newcommand{\pol}{\mathbb{P}[x]}                           				        % Polinomios
  \newcommand{\funl}{\mathcal{L}}                            				        % Funcional de momentos L
  \newcommand{\spol}{\{P_n\}_{n \in \N_0}}                   				        % Sucesión de polinomios
  \newcommand{\pescalar}[2]{\langle #1, #2 \rangle}      				            % Producto escalar
  \newcommand{\charlier}{\{C_n^{(\mu)}\}_{n \in \N_0}}       			   	      % SPOVD de Charlier
  \newcommand{\meixner}{\{M_n^{(\beta, c)}\}_{n \in \N_0}}                 	% SPOVD de Meixner
  \newcommand{\kraw}{\{K_n^{(p)}\}_{n \in \{0, \dots, N-1\}}}              	% SPOVD de Krawtchouk
  \newcommand{\hahn}{\{h_n^{(\alpha, \beta)}\}_{n \in \{0, \dots, N-1\}}}	  % SPOVD de Hahn
  \newcommand{\eprob}{(\Omega, \mathcal A, P)}              				        % Espacio de probabilidad
  \newcommand{\pest}{\{X_t\}_{t \in T}}                     				        % Proceso estocástico
  \newcommand{\markov}{\{X_t\}_{t \in \mathbb{N}_0}}        				        % Cadena de Markov
  \newcommand{\pnm}{\{X_t\}_{t \in [0, \infty[}}            				        % Proceso de nacimiento y muerte

% Para escalar matemáticas:
  \newcommand{\scalemath}[2]{\scalebox{#1}{\mbox{\ensuremath{\displaystyle #2}}}}

% TEOREMAS Y ENTORNOS ASOCIADOS

  % \newtheorem{theorem}{Theorem}[chapter]
  \newtheorem*{teorema*}{Teorema}
  \newtheorem{teorema}{Teorema}[chapter]
  \newtheorem{proposicion}[teorema]{Proposición}
  \newtheorem{lema}[teorema]{Lema}
  \newtheorem{corolario}[teorema]{Corolario}

    \theoremstyle{definition}
  \newtheorem{definicion}[teorema]{Definición}
  \newtheorem{ejemplo}[teorema]{Ejemplo}

    \theoremstyle{remark}
  \newtheorem{observacion}[teorema]{Observación}


% --------------------------------------------------------------------
% INFORMACIÓN DEL TFG Y EL AUTOR
% --------------------------------------------------------------------
\usepackage{xspace} % Para problemas de espaciado al definir comandos

\newcommand{\miTitulo}{Polinomios Ortogonales y \\ Procesos de Nacimiento y Muerte. \\ Implementación en Python\xspace}
\newcommand{\miNombre}{José Luis Ruiz Benito\xspace}
\newcommand{\miGrado}{Doble Grado en Ingeniería Informática y Matemáticas}
\newcommand{\miFacultad}{Facultad de Ciencias y Escuela Técnica Superior de Ingeniería Informática y de Telecomunicación}
\newcommand{\miUniversidad}{Universidad de Granada}
% Añadir tantos tutores como sea necesario separando cada uno de ellos
% mediante el comando `\medskip` y una línea en blanco
\newcommand{\miTutor}{
  Teresa E. Pérez Fernández \\ \emph{Departamento de Matemática Aplicada} 

  % Añadir tantos tutores como sea necesario. 

  \medskip
  Francisco Luque Sánchez \\ \emph{Departamento de Ciencias de la \\ Computación e Inteligencia Artificial}
}
\newcommand{\miCurso}{2021-2022\xspace}

% HYPERREFERENCES
% --------------------------------------------------------------------
\usepackage{xurl}
\usepackage[pagebackref]{hyperref}
% Opciones para el paquete hyperref
%----------------------------------

\hypersetup{%
    % hidelinks,            % Enlaces sin color ni borde. El borde no se imprime
    linkbordercolor=0.8 0 0,
    citebordercolor=0 0.8 0,
    citebordercolor=0 0.8 0,
    colorlinks = true,            % Color en texto de los enlaces. Comentar esta línea o cambiar `true` por `false` para imprimir el documento.
    linkcolor = [rgb]{0.5, 0, 0}, % Color de los enlaces internos
    urlcolor = [rgb]{0, 0, 0.5},  % Color de los hipervínculos
    citecolor = [rgb]{0, 0.5, 0}, % Color de las referencias bibliográficas
    pdftitle={\miTitulo},%
    pdfauthor={\textcopyright\ \miNombre, \miFacultad, \miUniversidad},%
    pdfsubject={Trabajo de fin de Grado},%
    pdfkeywords={},%
    pdfcreator={pdfLaTeX},%
}

% Redefinición del estilo para mostrar las referencias cruzadas en la bibliografía.
\renewcommand*{\backref}[1]{}
\renewcommand*{\backrefalt}[4]{{\footnotesize [%
    \ifcase #1 No citado%
    \or Citado en pág.~#2%
    \else Citado en págs. #2%
    \fi%
]}}

% Etiquetas en español para el comando \autoref
\def\chapterautorefname{Capítulo}
\def\appendixautorefname{Apéndice}
\def\sectionautorefname{Sección}
\def\subsectionautorefname{Subsección}
\def\figureautorefname{Figura}
\def\tableautorefname{Tabla}

\def\teoremaautorefname{Teorema}
\def\proposicionautorefname{Proposición}
\def\lemaautorefname{Lema}
\def\corolarioautorefname{Corolario}
\def\definicionautorefname{Definición}
\def\observacionautorefname{Observación}
\def\ejemploautorefname{Ejemplo}

% Pone automáticamente un parántesis para las ecuaciones
\def\equationautorefname~#1\null{Ecuación ~(#1)\null}


\begin{document}

% --------------------------------------------------------------------
% FRONTMATTER
% -------------------------------------------------------------------
%\frontmatter % Desactiva la numeración de capítulos y usa numeración romana para las páginas

% \pagestyle{plain} % No imprime cabeceras

% !TeX root = ../libro.tex
% !TeX encoding = utf8

%*******************************************************
% Titlepage
%*******************************************************
\begin{titlepage}
  \AddToShipoutPicture*{\BackgroundPic}
  \phantomsection 
  \pdfbookmark[1]{Título}{title}

  % Para que el título esté centrado en la página.
  % Los valores numéricos deberán elegirse de acuerdo con el diseño de
  % página (sobre todo si se cambia la opción BCOR o DIV).
  \begin{addmargin}[2.575cm]{0cm}
  \begin{flushleft}
    \Large  
    \hfill\vfil

    \large{\textsf{\miFacultad}}
    \vfill

    {\large\textsc\miGrado} \vfill


    {\large\textsc{trabajo de fin de grado}}

    \begin{flushleft}
      \Huge
      \setstretch{0.8}
      \miTitulo
    \end{flushleft}

    \vfill\vfill\vfill\vfill

    \textsf{\normalsize{Presentado por:}}\\
    {\normalsize\textrm{\miNombre}} 
    \bigskip

    \textsf{\normalsize{Tutores:}}\\
    {\normalsize\rmfamily\miTutor}

    \bigskip
    \textsf{\normalsize{Curso académico \miCurso}}
  \end{flushleft}  
  \end{addmargin}       

\end{titlepage}   
\cleardoublepage
\endinput
           
% !TeX root = ../libro.tex
% !TeX encoding = utf8

%*******************************************************
% Little Dirty Titlepage
%*******************************************************

\thispagestyle{empty}

\begin{center}
  \large  

  \vspace*{\stretch{1}}

  \begingroup
  \huge{\miTitulo} \\ \bigskip
  \endgroup

  \textrm{\miNombre}

  \vspace{\stretch{5}}

\end{center}  

\newpage
\thispagestyle{empty}

\hfill

\vfill

\miNombre \textit{\miTitulo}.

Trabajo de fin de Grado. Curso académico \miCurso.
\bigskip

\begin{minipage}[t]{0.25\textwidth}
  \flushleft
  \textbf{Responsables de tutorización}
\end{minipage}
\begin{minipage}[t]{0.45\textwidth}
  \flushleft
  \miTutor
\end{minipage}
\begin{minipage}[t]{0.30\textwidth}
  \flushright
  \miGrado
  \medskip

  \miFacultad
  \medskip

  \miUniversidad
\end{minipage}

\newpage
\endinput
                     
% !TeX root = ../libro.tex
% !TeX encoding = utf8
%
%*******************************************************
% Declaración de originalidad
%*******************************************************

\thispagestyle{empty}

\hfill\vfill

\textsc{Declaración de originalidad}\\\bigskip

D. \miNombre \\\medskip

Declara explícitamente que el trabajo presentado como Trabajo de Fin de Grado (TFG), correspondiente al curso académico \miCurso, es original, en el sentido de que no ha utilizado para la elaboración del trabajo fuentes sin citarlas debidamente.
\medskip

En Granada a \today 
\begin{flushleft} 
Fdo: \miNombre 

\end{flushleft}

\vfill

\cleardoublepage
\endinput
   
%% !TeX root = ../libro.tex
% !TeX encoding = utf8

%*******************************************************
% Dedication
%*******************************************************
\thispagestyle{empty}
\phantomsection 
\pdfbookmark[1]{Dedicatoria}{Dedicatoria}

\hfill
\vfill

\begin{flushright}

\itshape
A mi familia, a mis amigos \\
y a quienes no puedieron verme graduarme.

\end{flushright}

\vfill

\cleardoublepage
\endinput
                % Opcional
% !TeX root = ../libro.tex
% !TeX encoding = utf8

%*******************************************************
% Table of Contents
%*******************************************************
\phantomsection
\pdfbookmark[0]{\contentsname}{toc}

\setcounter{tocdepth}{2} % <-- 2 includes up to subsections in the ToC
\setcounter{secnumdepth}{3} % <-- 3 numbers up to subsubsections

% \manualmark
% \markboth{\textsc{\contentsname}}{\textsc{\contentsname}}
\tableofcontents 

%*******************************************************
% List of Figures and of the Tables
%*******************************************************

    % *******************************************************
    %  List of Figures
    % *******************************************************    
    \phantomsection 
    % \listoffigures

    %*******************************************************
    % List of Tables
    %*******************************************************
    \phantomsection 
    % \listoftables
    
    %*******************************************************
    % List of Listings
    % The package \usepackage{listings} is needed
    %*******************************************************      
	  % \phantomsection 
    % \renewcommand{\lstlistlistingname}{Listados de código}
    % \lstlistoflistings 

\cleardoublepage
            
%% !TeX root = ../libro.tex
% !TeX encoding = utf8

%*******************************************************
% Agradecimientos
%*******************************************************

\chapter{Agradecimientos}

Agradecimientos del libro (opcional, ver archivo \texttt{preliminares/agradecimiento.tex}).

\cleardoublepage
\endinput
            % Opcional

% \pagestyle{scrheadings} % A partir de ahora sí imprime cabeceras

% !TeX root = ../libro.tex
% !TeX encoding = utf8
%
%*******************************************************
% Summary
%*******************************************************

\selectlanguage{english}
\addchap{Summary}

Summary en preliminares/summary.tex

\selectlanguage{spanish} 
\endinput
                    
% !TeX root = ../libro.tex
% !TeX encoding = utf8
%
%*******************************************************
% Introducción
%*******************************************************
% \manualmark
% \markboth{\textsc{Introducción}}{\textsc{Introducción}} 
\addchap{Introducción}

Introducción en preliminares/introduccion.tex

\endinput
               

% --------------------------------------------------------------------
% MAINMATTER
% --------------------------------------------------------------------
%\mainmatter % activa la numeración de capítulos, resetea la numeración de las páginas y usa números arábigos

%\setpartpreamble[c][0.75\linewidth]{%
%	\bigskip % Deja un espacio vertical en la parte superior
%  Si el trabajo se divide en diferentes partes es posible incluir al inicio de cada una de ellas un breve resumen que indique el contenido de la misma. Esto es opcional.
%}

\part{Polinomios Ortogonales y \\ Procesos de Nacimiento y Muerte} % Dividir un libro en partes OPCIONAL
\cleardoublepage    
\input{capitulos/matematicas/polinomios_ortogonales}
\input{capitulos/matematicas/variable_discreta}
\input{capitulos/matematicas/procesos}

\cleardoublepage\part{Implementación de una biblioteca de Polinomios Ortogonales de Variable Discreta en Python}
\chapter{Análisis del problema}\label{ch:cuarto-capitulo}
    
    Capítulo en memoria/capitulos/analisis.tex

\endinput
\input{capitulos/informatica/solucion}
\input{capitulos/informatica/demostracion}


% Conclusiones
\addchap{Conclusiones}

    Durante el desarrollo de este trabajo, hemos realizado un análisis de los fundamentos teóricos y prácticos de algunos criptosistemas de clave pública, poniendo especial atención al criptosistema de Merkle-Hellman. Dado que este criptosistema se basa en el problema de la mochila, hemos abarcado una gran variedad de conceptos necesarios para su compresión, que van desde nociones como las máquinas de Turing, hasta cuestiones tan fundamentales como el problema P vs NP.

    A lo largo de este proyecto, además de explicar toda la teoría matemática apoyada de ejemplos esclarecedores, hemos destacado la relevancia práctica de los criptosistemas mediante su implementación. Esto nos ha permitido analizar los resultados obtenidos y compararlos con los teóricos esperados. De esta manera, hemos conseguido plasmar todo el conocimiento estudiado en aplicaciones concretas.

    Asimismo, hemos abordado de manera firme los ataques principales a los criptosistemas estudiados, entendiendo que resulta casi imposible garantizar la total seguridad de un sistema criptográfico. Incluso en la actualidad, persistimos en la búsqueda de dicho criptosistema, especialmente ante el desafío que representan los algoritmos cuánticos, los cuales ponen en riesgo a muchos de los criptosistemas actuales.
    
    Para finalizar, este trabajo de fin de grado ha supuesto el primer acercamiento a la actividad profesional, y como tal, este proyecto queda abierto ante futuras investigaciones y desarrollos en el campo de la criptografía y de la computación cuántica.
    
\endinput



% --------------------------------------------------------------------
% APPENDIX: Opcional
% --------------------------------------------------------------------

\appendix % Reinicia la numeración de los capítulos y usa letras para numerarlos
\pdfbookmark[-1]{Apéndices}{appendix} % Alternativamente podemos agrupar los apéndices con un nuevo \part{Apéndices}

%\input{apendices/apendice01}
% Añadir tantos apéndices como sea necesario 
%\input{apendices/codigo}

% --------------------------------------------------------------------
% GLOSARIO: Opcional
% --------------------------------------------------------------------

%\input{glosario} 

% -------------------------------------------------------------------
% BACKMATTER
% -------------------------------------------------------------------

\backmatter % Desactiva la numeración de los capítulos
\pdfbookmark[-1]{Referencias e Índices}{BM-Referencias}

% BIBLIOGRAFÍA
%-------------------------------------------------------------------

\setbibpreamble{Las referencias se listan por orden alfabético. Aquellas referencias con más de un autor están ordenadas de acuerdo con el primer autor.\par\bigskip}
\bibliographystyle{alpha}
\begin{small} % Normalmente la bibliografía se imprime en un tamaño de letra más pequeño.
  \bibliography{library.bib}
\end{small}


% ÍNDICE TERMINOLÓGICO  (Opcional) 
%------------------------------------------------------------------- 

% Para incluir el índice terminológico es necesario descomentar los siguientes comandos. Incluir un índice terminológico es opcional

% \cleardoublepage 
% \begin{footnotesize} % Normalmente el índice se imprime en un tamaño de letra más pequeño.
% \printindex 
% \end{footnotesize}

\end{document}
